\chapter{The Locking Crisis}\label{ch:locking-crisis}

This chapter makes the software-engineering case against lock-based concurrency
and for declarative atomicity. The argument is not that locks are useless;
they are precise and well understood. The argument is that locks force the
programmer to maintain a synchronization protocol that is invisible to the
compiler, the type system, and every future reader of the code. By
``invisible'' we mean three concrete things. The compiler cannot see the
protocol: a \texttt{lock()}/\texttt{unlock()} pair is, to it, just two
independent function calls, so it has no basis to reorder, cache, or eliminate
accesses around them, and no way to warn that a field is being touched without
its lock. The type system cannot express it: nothing in a field's type says
``this member is protected by \texttt{g\_lock}, and only while it is held,''
so the compiler cannot reject an unsynchronized access the way it rejects a
type mismatch. And the protocol is not recoverable by inspection: the lock
order, the set of guarded fields, and the conditions under which threads may
wait are spread across many functions and modules, so a reader must reconstruct
them from scattered acquire and release sites, and any two readers may
reconstruct different protocols. When synchronization is invisible in this
way, the language and its tools cannot help; every invariant must be enforced
by convention alone \cite{boehm05}. The moment two separately correct lock-based
modules interact, their lock orders can compose into a deadlock that no local
inspection of either module can prevent. The rest of this book develops an
alternative: transactional memory \cite{herlihy93}, in which the programmer
declares the \emph{set} of accesses that must appear atomic rather than an
\emph{order} among locks.

\section{Why Synchronization Became a First-Class Problem}
\index{history of concurrency}
\index{distributed shared memory}
\index{cache coherence!history}
\index{multicore!history}

Shared memory is so familiar today that it is easy to forget that it had to be
invented, and that the invention took several decades. The path runs through
three eras, and each one left its stamp on how we synchronize. In the
\textbf{distributed era}, processors lived in separate machines and shared
nothing but a network; coordination meant explicit message passing, and the
central problem was to make a set of machines with independent clocks and no
common memory agree on anything at all \cite{lamport78}. Distributed systems
produced the first general theories of ordering and agreement, but they gave
programmers no shared address space to think in.

The \textbf{distributed shared-memory era} was an attempt to fix exactly that.
Researchers built systems that \emph{simulated} a single address space across
a network, using software protocols --- ownership transfer, invalidation,
release consistency --- to keep the simulation coherent
\cite{li89dsm,amza96treadmarks}. These systems made the shared-memory
programming model available on hardware that did not actually share memory, at
the cost of high communication overhead. They were a programming convenience,
not a hardware change.

The \textbf{multicore era} changed the hardware instead. When it became
cheaper to put several processors on one chip than to make a single processor
faster \cite{olukotun96case}, the industry built multicore chips in which
cores \emph{physically} share memory through a hardware cache-coherence
protocol \cite{goodman83,papamarcos84,archibald86}. Shared memory stopped
being a simulation and became the machine itself. This is the moment that
synchronization stopped being a niche concern of distributed-systems
researchers and became a first-class problem for every programmer: thousands
of cores, one address space, and only locks --- and, as we will see, their
successors --- to keep the accesses coherent. The rest of this book is an
extended answer to the question that this era posed. Chapter~2 recounts the
full history in more detail; here it is enough to note that the hardware now
gives us shared memory for free, and the remaining difficulty is purely in how
we program it.

\section{The Cost of Mixing Synchronization with Logic}
\index{locking crisis}
\index{deadlock}
\index{priority inversion}
\index{lost wake-up}
\index{data race}
\index{Heisenbug}

A lock-based critical section is a small program written in a second language:
a language of acquire and release, wait and notify, trylock and unlock. The
programmer must speak both languages at once. The application logic says what
the program should compute; the synchronization logic says which interleavings
are legal. The two are entangled in the same source text, and that entanglement
is the root of the software-engineering problem.

Real-world systems fail in every one of these ways. Priority inversion stalled
the Mars Pathfinder probe for days in 1997 because a low-priority task held a
mutex needed by a high-priority task \cite{jones97mars,sha90priority}. Lost
wake-ups corrupt sleeping synchronization when a signal arrives before the
waiting thread blocks. Unrelated defects in production concurrency code
routinely cost organizations weeks of debugging time precisely because the
failure does not leave an obvious trail. None of these incidents reflects an
exotic language feature; each is a normal consequence of entangling logic with
an invisible protocol.

Three failure modes matter in practice. \textbf{Deadlock} means no progress at
all: two or more threads wait for conditions that can never become true, and
it is often the easiest failure to detect because the program simply stops.
\textbf{Livelock and starvation} mean progress exists but is not fair: a
thread may make no progress while other threads proceed indefinitely, as when a
spinlock owner is preempted while other threads spin, or in a
priority-inversion chain in which a low-priority thread holds a lock needed by
a high-priority thread. \textbf{Data races} mean incorrectness: a race occurs
when two threads access the same memory without ordering, at least one writes,
and the result depends on the interleaving. A race is not merely a theoretical
hazard; it is undefined behavior in most modern languages
\cite{boehm08memory}.

In each case, the synchronization contract is invisible. The compiler cannot
tell which lock guards which variable. The type system cannot express that a
field may be accessed only while holding a particular mutex. The reader cannot
recover the intended lock order from the code because the order is distributed
across many functions and modules.

Several error-prone patterns recur. In \textbf{check-then-act}, a thread reads
a value, decides, and then writes; the window between read and write is a race
unless the whole sequence is made atomic. In \textbf{test-and-set loops}, a
thread repeatedly tests a condition and attempts a compare-and-swap; the loop
is easy to write incorrectly and can suffer from starvation. \textbf{Nested
locking} arises when a function that acquires a lock then calls another
function that acquires another lock, creating an implicit lock-order
obligation that is not visible at either call site. In \textbf{lock-order
inversion}, two code paths acquire the same locks in different orders; each
path is locally correct, but their composition can deadlock.

Concurrency bugs reproduce only under specific interleavings. This makes them
\emph{Heisenbugs}: changing the timing, adding a log statement, or running
under a debugger can make the bug disappear. The cost of such bugs is
disproportionately high because they survive review, unit tests, and even
integration tests, only to appear in production.

There is also a hidden performance cost. Coarse-grained locks serialize
executions and reduce available parallelism. The expected speedup of a program
with a serial fraction $f$ on $N$ cores is bounded by Amdahl's law
\cite{amdahl67}, discussed later in Part~IV. The cost of lock contention is
therefore not only a software engineering concern; it is a performance concern
as well.

\section{Composability -- The ``Innocent Change'' Problem}
\index{composability}
\index{lock order inversion}
\index{nested locking}

Composability is the ability to combine two correct components into a correct
larger component without modifying either one. Lock-based code is not
composable. Two separately correct critical sections can deadlock when
composed because each critical section imposes a lock-order obligation that is
not part of its public interface.

Consider a thread-safe \texttt{List} and a thread-safe \texttt{Map}. Each is
correct in isolation: the \texttt{List} uses one lock, the \texttt{Map} uses
another. Now suppose a program must move an element from the \texttt{List} to
the \texttt{Map}. The operation must lock both. But which lock first? If the
programmer chooses \texttt{List} first and \texttt{Map} second, then every
other path that touches both structures must use the same order. If another
path uses \texttt{Map} first, the two paths can deadlock. The correctness of
the composed operation depends on a global law that is not enforced by either
component and is not visible to the compiler.

The decisive advantage of transactional memory is that it removes the
lock-order obligation. A transaction names the set of memory locations it
accesses. Conflict is detected on the accessed data at runtime, not by an
application-wide order among locks. Two transactions that access disjoint data
do not block each other. Two transactions that access the same data are
serialized by the runtime according to a well-defined conflict protocol.

Nesting should be free. An atomic block inside an atomic block should flatten
into one atomic block. With locks, nested critical sections must be carefully
avoided or carefully ordered. With transactions, the composition rule is
simple: the outer transaction is the only real boundary, and the inner
transaction is an ordinary function call that inherits atomicity from its
caller.

\section{Running Example: The Innocent Change in Code}
\index{bank transfer}

The rest of this book carries one small example through every chapter: a
shared bank account transfer. The initial version is intentionally simple:

\begin{lstlisting}[language=C, caption={The global-lock transfer.}, label={lst:transfer-basic}]
void transfer(struct account *a, struct account *b, int n) {
    lock(&g_lock);
    a->balance -= n;
    b->balance += n;
    unlock(&g_lock);
}
\end{lstlisting}

The code is correct as long as every transfer uses the single global lock. Now
an ``innocent'' change adds a business rule: a transfer must not push either
account below a floor. The programmer adds a second lock to protect the floor
state:

\begin{lstlisting}[language=C, caption={The innocent change: two locks, an ordering obligation.}, label={lst:transfer-safe}]
void transfer_safe(struct account *a, struct account *b, int n) {
    lock(&g_lock);
    lock(&g_floor_lock);
    if (a->balance >= floor + n && b->balance >= floor) {
        a->balance -= n;
        b->balance += n;
    }
    unlock(&g_floor_lock);
    unlock(&g_lock);
}
\end{lstlisting}

Somewhere else, an audit path checks the floor before checking balances and
locks the two locks in the opposite order. The two paths together create a
wait-for cycle, shown in Figure~\ref{fig:two-lock-deadlock}.

\begin{figure}[htbp]
\centering
\begin{tikzpicture}[node distance=2.4cm, >=Latex]
\node[draw, rounded corners, align=center, minimum height=1.2cm] (transfer)
at (0,0) {Thread $T_1$ holds \texttt{g\_lock},\\requests \texttt{g\_floor\_lock}};
\node[draw, rounded corners, align=center, minimum height=1.2cm] (audit)
at (7,0) {Thread $T_2$ holds \texttt{g\_floor\_lock},\\requests \texttt{g\_lock}};
\draw[->, bend left=20] (transfer) to node[above,midway,font=\small]
{waits for \texttt{g\_floor\_lock}} (audit);
\draw[->, bend left=20] (audit) to node[below,midway,font=\small]
{waits for \texttt{g\_lock}} (transfer);
\end{tikzpicture}
\caption{The two-lock wait-for cycle created by the innocent change. Neither
critical section is wrong by itself; the deadlock emerges from the
composition of two opposite lock orders.}
\label{fig:two-lock-deadlock}
\end{figure}

The same business logic can be expressed as a transaction:

\begin{lstlisting}[language=C, caption={The same logic as a transaction.}, label={lst:transfer-atomic}]
void transfer_atomic(struct account *a, struct account *b, int n) {
    __transaction_atomic {
        if (a->balance >= floor + n && b->balance >= floor) {
            a->balance -= n;
            b->balance += n;
        }
    }
}
\end{lstlisting}

The transaction names the set of accessed addresses. The runtime detects
conflicts on the data actually touched; the lock-order obligation is gone.
Table~\ref{tab:locking-crisis-taxonomy} summarizes the shift in burden.

\begin{table}[htbp]
\centering
\caption{The locking-crisis taxonomy: each synchronization approach shifts
a different burden onto the programmer, the compiler, or the runtime.}
\label{tab:locking-crisis-taxonomy}
\footnotesize
\renewcommand{\arraystretch}{1.35}
\setlength{\tabcolsep}{4pt}
\begin{tabular}{M{0.13\linewidth} M{0.18\linewidth} M{0.14\linewidth} M{0.16\linewidth} M{0.13\linewidth}}
\hline
\thb{0.13\textwidth}{\textbf{Approach}} & \thb{0.18\textwidth}{\textbf{Synchronization order}} & \thb{0.14\textwidth}{\textbf{Deadlock risk}} & \thb{0.16\textwidth}{\textbf{Composability}} & \thb{0.13\textwidth}{\textbf{Reasoning locality}} \\
\hline
Single global lock & none & none & coarse but safe & global \\
Ordered multiple locks & programmer-defined & high if order violated & not composable & distributed \\
Transactional \texttt{atomic} & runtime-detected & none & composable & local \\
\hline
\end{tabular}
\end{table}

\section{Worked Example: Reproducing the Two-Lock Deadlock}
\index{two-lock deadlock}

The following program is deliberately small and compilable. It demonstrates
that the deadlock is not caused by either critical section alone, but by the
opposite lock orders of \texttt{transfer\_safe} and \texttt{audit}.

\begin{lstlisting}[language=C, caption={A compilable two-lock deadlock. Build with \texttt{cc -std=c11 -pthread deadlock.c -o deadlock}.}, label={lst:two-lock-deadlock}]
#include <pthread.h>
#include <stdio.h>
#include <unistd.h>

pthread_mutex_t g_lock = PTHREAD_MUTEX_INITIALIZER;
pthread_mutex_t g_floor_lock = PTHREAD_MUTEX_INITIALIZER;

int balance_a = 100;
int balance_b = 0;
int floor = 0;

void *transfer_safe(void *arg) {
    int n = *(int *) arg;
    pthread_mutex_lock(&g_lock);
    usleep(1000);
    pthread_mutex_lock(&g_floor_lock);
    if (balance_a >= floor + n) {
        balance_a -= n;
        balance_b += n;
    }
    pthread_mutex_unlock(&g_floor_lock);
    pthread_mutex_unlock(&g_lock);
    return NULL;
}

void *audit(void *arg) {
    pthread_mutex_lock(&g_floor_lock);
    usleep(1000);
    pthread_mutex_lock(&g_lock);
    int total = balance_a + balance_b;
    pthread_mutex_unlock(&g_lock);
    pthread_mutex_unlock(&g_floor_lock);
    printf("total = %d\n", total);
    return NULL;
}

int main(void) {
    pthread_t t1, t2;
    int amount = 10;
    pthread_create(&t1, NULL, transfer_safe, &amount);
    pthread_create(&t2, NULL, audit, NULL);
    pthread_join(t1, NULL);
    pthread_join(t2, NULL);
    return 0;
}
\end{lstlisting}

The \texttt{usleep} calls do not create the bug. They only widen the bad
interleaving window so that the deadlock is easier to observe. A global lock
order can repair this program, but that order becomes an application-wide law:
every future function that touches both locks must obey it, and the law is not
checked by the compiler. The transactional memory designs in later chapters
abandon the lock-order law altogether.

\section{Worked Example: Model Checking the Bank Invariant}
\index{model checking}
\index{TLA+}

The locking bug is an implementation artifact. The intended application
invariant is simple: an accepted transfer preserves total money and never
leaves an account below the floor. A model of the atomic transfer states
exactly that, without mutexes and without a lock order.

\begin{lstlisting}[caption={A TLA+ model of atomic bank transfers and their invariants.}, label={lst:bank-atomic}]
---- MODULE BankAtomic ----
EXTENDS Integers, TLC

CONSTANTS InitBal, Floor, Amt

VARIABLES balA, balB
vars == <<balA, balB>>

Init ==
/\ balA = InitBal
/\ balB = InitBal

Transfer ==
/\ balA >= Floor + Amt
/\ balA' = balA - Amt
/\ balB' = balB + Amt

Next ==
/ Transfer
/ UNCHANGED vars

MoneyInvariant ==
balA + balB = 2 * InitBal

FloorInvariant ==
/\ balA >= Floor
/\ balB >= Floor

Spec == Init /\ [][Next]_vars
====
\end{lstlisting}

The model states the semantic invariant directly. The \texttt{Transfer} action
corresponds to the source-level transaction, not to a particular TM algorithm.
Later chapters refine this step into concrete mechanisms: SGL, NOrec, TL2,
TinySTM, SwissTM, and the other backends all implement different
conflict-detection and commit protocols, but every one of them must preserve
the same two invariants. Verification should separate the application
invariant from the synchronization mechanism whenever possible.

\section{Summary}
\label{sec:ch1-summary}

Lock-based synchronization entangles concurrency control with application
logic. That entanglement produces three failure modes: deadlock, starvation,
and data races. The deepest problem is not that any single lock order is hard
to write; it is that lock-based code is not composable. Two locally correct
modules can compose into a deadlocking program, and no amount of local
reasoning prevents it.

Transactional memory changes the programming model. A transaction names the
set of accesses that must appear atomic. The runtime detects conflicts on the
accessed data, removing the lock-order obligation. The same transfer logic
that required a global lock order becomes a declarative atomic block, and
every Part~IV runtime implements exactly that small API. The rest of this book
makes the intuition precise and then replaces it with mechanisms: compilers
that instrument loads and stores, runtimes that manage speculative state, and
formal models that prove the application invariant survives concurrency.

\section{Exercises}
\label{sec:ch1-exercises}

\begin{enumerate}
\item \emph{[theory]} Describe a scenario in which two independently
correct lock-based modules become deadlocked when composed. Explain
why the deadlock cannot be found by testing either module alone.

\item \emph{[theory]} ``Just use finer-grained locks.'' Discuss the limits
of this advice. Give one concrete example where finer-grained locking
reduces contention but increases the probability of lock-order
inversion.

\item \emph{[implementation]} Compile the program in
Listing~\ref{lst:two-lock-deadlock}. Replace the opposite-order audit
path with a single global lock order, then with a
\texttt{pthread\_mutex\_trylock} rollback path. Explain which version
removes deadlock and which version can still starve.

\item \emph{[verification]} Extend the TLA+ model in
Listing~\ref{lst:bank-atomic} with a rejected-transfer action that
leaves \texttt{balA} and \texttt{balB} unchanged when the floor check
fails. Model check that \texttt{MoneyInvariant} and
\texttt{FloorInvariant} remain inductive.

\item \emph{[theory]} Give a precise argument for why nested locks do not
compose automatically, while nested atomic blocks can be defined by
flattening. State one assumption about the TM runtime that this
argument requires.

\item \emph{[implementation]} Run the two-lock deadlock program under the
\texttt{rr} record-and-replay debugger. Identify the program point at
which each thread blocks, and explain how a wait-for-cycle graph can be
computed from a stack trace.

\item \emph{[verification]} The TLA+ model in Listing~\ref{lst:bank-atomic}
splits a transfer into two visible steps if you remove
\texttt{Transfer} and expose two separate actions. Show that TLC finds
a state that violates \texttt{MoneyInvariant}. Explain which
source-level instrumented read/write pair the two actions represent.

\item \emph{[open]} The ``locking crisis'' is usually described as a
failure of \emph{composability}.  Argue a position: does transactional
memory \emph{solve} the crisis or merely \emph{move} it into the
runtime, where the remaining problems (reclamation, retries, abort
side effects, nesting) are simply harder to see?  State the assumption
about programmers or about runtimes your position rests on, and give
the strongest argument against it.
\end{enumerate}% -------------------------------------------------------------------

\textbf{How to check your work.} The \emph{[implementation]} exercises are
verified with the companion: compile Listing~\ref{lst:two-lock-deadlock} and
run \texttt{make -C tests run}.  The \emph{[verification]} exercises are
checked with TLC against Listing~\ref{lst:bank-atomic}
(\texttt{MoneyInvariant}, \texttt{FloorInvariant}; see
Appendix~\ref{chap:pluscal-quick-ref}).  The \emph{[theory]} exercises are
graded against \S\ref{ch:locking-crisis}.  The \emph{[open]} exercise has
no single right answer; it is graded on whether you state the assumption
your position rests on and address the strongest counterargument.

\index{composability}
